\documentclass[preview]{standalone}

\usepackage{amsmath}
\usepackage{amssymb}
\usepackage{stellar}
\usepackage{definitions}

\begin{document}

\id{module-submodules-morphisms-quotients}
\genpage

\section{Submodules}

\begin{snippetdefinition}{submodule-definition}{Submodule}
    Let \(A\) be a \ring and let \(M\) be a left \(A\)-\module. A
    \set \(N\subseteq M\) is an \emph{\(A\)-submodule} of \(M\) if
    \(N\) is an additive \subgroup of \(M\) and
    \[
        av\in N
        \qquad
        \text{for every }a\in A\text{ and }v\in N.
    \]
    In this case one writes
    \[
        N\submoduleleq M.
    \]
\end{snippetdefinition}

\begin{snippetdefinition}{module-homomorphism-definition}{Module homomorphism}
    Let \(A\) be a \ring and let \(M_1,M_2\) be left
    \(A\)-\module[modules]. A
    \function
    \[
        \varphi\colon M_1\fromto M_2
    \]
    is an \emph{\(A\)-module homomorphism} if, for every
    \(v,w\in M_1\) and \(a\in A\),
    \[
        \varphi(v+w)=\varphi(v)+\varphi(w),
        \qquad
        \varphi(av)=a\varphi(v).
    \]
    Its kernel and image are
    \[
        \moduleker\varphi
        =
        \{v\in M_1\suchthat\varphi(v)=0_{M_2}\},
        \qquad
        \moduleimage\varphi
        =
        \{\varphi(v)\suchthat v\in M_1\}.
    \]
    A bijective module homomorphism is a \emph{module isomorphism}.
\end{snippetdefinition}

\section{Quotient modules}

\begin{snippetdefinition}{quotient-module-definition}{Quotient module}
    Let \(A\) be a \ring, let \(M\) be a left \(A\)-\module, and let
    \(N\submoduleleq M\). The quotient additive group \(M/N\) becomes a
    left \(A\)-\module under
    \[
        a\cdot(v+N)=av+N.
    \]
    It is called the \emph{quotient module} of \(M\) by \(N\).
\end{snippetdefinition}

\begin{snippetproof}{quotient-module-definition-proof}{quotient-module-definition}{Quotient module}
    If \(v+N=v'+N\), then \(v-v'\in N\). Since \(N\) is an
    \(A\)-\submodule,
    \[
        a(v-v')=av-av'\in N.
    \]
    Therefore \(av+N=av'+N\), so scalar multiplication is well-defined.
    The module axioms follow from those of \(M\).
\end{snippetproof}

\section{Operations and isomorphism theorems}

\begin{snippetproposition}{submodule-operations-and-homomorphism-kernels}{Operations with submodules}
    Let \(A\) be a \ring and let \(M,M'\) be left
    \(A\)-\module[modules].
    \begin{enumerate}
        \item If \(N_1,N_2\submoduleleq M\), then
        \[
            N_1\intersection N_2\submoduleleq M
        \]
        and
        \[
            N_1+N_2
            =
            \{v_1+v_2\suchthat v_1\in N_1,\ v_2\in N_2\}
            \submoduleleq M.
        \]
        \item If \(\varphi\colon M\fromto M'\) is an
        \(A\)-\modulehomomorphism, then
        \[
            \moduleker\varphi\submoduleleq M,
            \qquad
            \moduleimage\varphi\submoduleleq M'.
        \]
        Moreover,
        \[
            \varphi\text{ is injective}
            \ifandonlyif
            \moduleker\varphi=\{0_M\}.
        \]
    \end{enumerate}
\end{snippetproposition}

\begin{snippetproof}{submodule-operations-and-homomorphism-kernels-proof}{submodule-operations-and-homomorphism-kernels}{Operations with submodules}
    The intersection is an additive subgroup stable under scalar
    multiplication. For the sum, if \(v_i\in N_1\), \(w_i\in N_2\), and
    \(a\in A\), then
    \[
        (v_1+w_1)-(v_2+w_2)
        =
        (v_1-v_2)+(w_1-w_2)\in N_1+N_2
    \]
    and
    \[
        a(v_1+w_1)=av_1+aw_1\in N_1+N_2.
    \]

    If \(v,w\in\moduleker\varphi\), then
    \[
        \varphi(v-w)=0_{M'},
        \qquad
        \varphi(av)=a\varphi(v)=0_{M'},
    \]
    so the kernel is a \submodule. If
    \(u_i=\varphi(v_i)\in\moduleimage\varphi\), then
    \[
        u_1-u_2=\varphi(v_1-v_2),
        \qquad
        au_1=\varphi(av_1),
    \]
    so the image is a \submodule. The kernel is zero exactly when
    \(\varphi(v)=\varphi(w)\) implies \(v=w\).
\end{snippetproof}

\begin{snippettheorem}{module-isomorphism-theorems}{Isomorphism theorems for modules}
    Let \(A\) be a \ring and let all \module[modules] below be left
    \(A\)-\module[modules].
    \begin{enumerate}
        \item If \(\varphi\colon M\fromto M'\) is a
        \modulehomomorphism,
        then
        \[
            M/\moduleker\varphi
            \moduleisomorphic
            \moduleimage\varphi.
        \]
        \item If \(N_1,N_2\submoduleleq M\), then
        \[
            \frac{N_1+N_2}{N_1}
            \moduleisomorphic
            \frac{N_2}{N_1\intersection N_2}.
        \]
        \item If \(N_1\submoduleleq N_2\submoduleleq M\), then
        \[
            \frac{M/N_1}{N_2/N_1}
            \moduleisomorphic
            M/N_2.
        \]
    \end{enumerate}
\end{snippettheorem}

\begin{snippetproof}{module-isomorphism-theorems-proof}{module-isomorphism-theorems}{Isomorphism theorems for modules}
    For the first statement, define
    \[
        \overline{\varphi}\colon
        M/\moduleker\varphi
        \fromto
        \moduleimage\varphi,
        \qquad
        \overline{\varphi}(v+\moduleker\varphi)=\varphi(v).
    \]
    It is a well-defined bijective \(A\)-\modulehomomorphism.

    For the second statement, the surjective homomorphism
    \[
        \psi\colon N_2\fromto(N_1+N_2)/N_1,
        \qquad
        \psi(v)=v+N_1,
    \]
    has kernel \(N_1\intersection N_2\). The first statement now gives
    the desired isomorphism.

    Finally, the surjective homomorphism
    \[
        M/N_1\fromto M/N_2,
        \qquad
        v+N_1\longmapsto v+N_2,
    \]
    has kernel \(N_2/N_1\). Applying the first statement once more gives
    the third isomorphism.
\end{snippetproof}

\end{document}
